\scanpage{46}
\mspage{2}

\uline{Objets}: Triples $(X,\widetilde\Sigma,\deleted{\unsure{G|S}}R)$, avec

\margin{$X,\deleted{\partial X}\!,\Sigma,S$}

$X$ surface à bord compacte orientée

$\widetilde\Sigma$ sous-$1$-variété compacte de $X$ \deleted{$\unsure{\partial X}$}, $\widetilde\Sigma\supset\partial X$

\deleted{$\unsure{G|S}$} $R\subset\Sigma$ \deleted{$\unsure{\partial X}$}, partie finie

% slanted
\margin{$\widetilde\Sigma=\partial X\amalg\Sigma$. On peut considérer que $R$ rigidifie, \unsure{mais} aussi que $R$ ``découpe'' $\widetilde\Sigma$\ldots}

\uline{flèches}: classes d'isotopie d'homéom. conservant l'orientation

\deleted{Variante} \deleted{\unsure{on}} Les \unsure{comp.}\,\unsure{des|s.} connexes de $\widetilde\Sigma\setminus\partial X$ \inserted{$=\Sigma$} (\inserted{celles de $\partial X$ sont \uline{les cercles de raccord}}) sont les \uline{cercles de découpage}, les éléments de \deleted{$\unsure{G|S}$}$R$ sont les \deleted{\unsure{points marqués}} \uline{sommets de} \deleted{\unsure{marquage}} \inserted{\uline{rigidification}} \deleted{(ou de rigidification)}, \deleted{les composantes de $\partial X$ sont les ``trous''}. Les composantes connexes de $X$ sont les \deleted{\unsure{morceaux}} \inserted{\uline{composants}} de l'objet envisagé.

Variante $S_{*}\unsure{\mathscr T|\mathbb T}{\unsure{\mathbb D|D}}_{\mathrm{top}}^{\pm}$, avec
\diagram{\begin{tikzcd}[ampersand replacement=\&,column sep=large]
 S_{*}\mathscr T\mathbb D_{\mathrm{top}}\arrow[r,hook,no head] \& S_{*}\mathscr T\mathbb D_{\mathrm{top}}^{\pm}\arrow[r] \& {[\pm1]}
\end{tikzcd}}
comme tantôt.

On a des sous-groupoïdes pleins

\diagram{\(\begin{array}{llll}
\deleted{\unsure{G}=\Sigma=\varnothing}\;\inserted{R=\varnothing,\Sigma=\varnothing}
 &\mathscr T_{\mathrm{top}} &\text{et } \mathscr T_{\mathrm{top}}^{\pm} &\\
\substack{\Sigma=\deleted{\unsure{\partial X}}\varnothing\\ R\text{ quelconque}}
 &S_{*}\mathscr T_{\mathrm{top}} &\text{et }S_{*}\mathscr T_{\mathrm{top}}^{\pm}
 &\substack{\text{(groupoïdes de Teichmüller }\text{\uline{spéciaux}}\\
 \text{\inserted{(ou \uline{rigidifiés})}}\\\text{(topologique orienté)}}\\
 \deleted{\unsure{S}}R=\varnothing
 &\mathscr T\mathbb D_{\mathrm{top}} &\text{et }\mathscr T\mathbb D_{\mathrm{top}}^{\pm}&
\end{array}\)}

(Cas pas intéressant je crois, sauf pour $\Sigma=\varnothing$ également)

% slanted
\margin{$MD={}$``\unsure{modulPair|modulaire}'' ou Mumford-Deligne}

\deleted{\unsure{Dé}} \uline{Détermination} du $\pi_0(S_{*}\mathscr T\mathbb D_{\mathrm{top}})\xrightarrow{\sim}\pi_0(S_{*}\mathscr T\mathbb D_{\mathrm{top}}^{\pm})$.

On introduit les \uline{graphes} \deleted{\unsure{MD}} bipondérés
% archive sheet number: 45

